\documentclass[a4paper]{article}     
\usepackage{amsfonts}
\usepackage{amsmath}
\usepackage{amssymb}

\newcommand{\adR}{\Rightarrow}
\newcommand{\Lh}{\left(}
\newcommand{\Rh}{\right)}
\newcommand{\intgrl}{\displaystyle\int}

\begin{document}

\noindent \large Auteur: Tala Mohanna \\
\noindent \large Studierichting: Informatica\\
\noindent \large Oefeningenreeks 4A nr. 38.3\\

\medskip

\normalsize

$ I = \intgrl \dfrac{2x+1}{4x^2-4x+3} dx$\\

$= \intgrl \dfrac{\dfrac{1}{4}\Lh 8x-4\Rh +2}{4x^2-4x+3} dx$\\

$= \dfrac{1}{4}\intgrl \dfrac{8x-4}{4x^2-4x+3} dx + 2\intgrl \dfrac{1}{4x^2-4x+3} dx$\\

Stel: $t = 4x^2-4x+3 \adR dt = \Lh 8x-4\Rh dx$\\

$= \dfrac{1}{4}\intgrl \dfrac{dt}{t} + 2\intgrl \dfrac{1}{4x^2-4x+3} dx$\\

$= \dfrac{1}{4}\ln |t| + 2\intgrl \dfrac{1}{\Lh 2x-1\Rh^2+2} dx$\\

$= \dfrac{1}{4}\ln |4x^2-4x+3| + 2\intgrl \dfrac{1}{\Lh 2x-1\Rh^2+2} dx$\\

Stel: $u = 2x-1 \adR du = 2dx \adR dx = \dfrac{du}{2}$\\

$= \dfrac{1}{4}\ln |4x^2-4x+3| + 2\intgrl \dfrac{1}{u^2+2}\cdot \dfrac{du}{2}$\\

$= \dfrac{1}{4}\ln |4x^2-4x+3| + \intgrl \dfrac{1}{u^2+\Lh \sqrt{2}\Rh^2}du$\\

$= \dfrac{1}{4}\ln |4x^2-4x+3| + \dfrac{1}{\sqrt{2}}\operatorname{Bgtan}\Lh \dfrac{u}{\sqrt{2}}\Rh + c$\\

$= \dfrac{1}{4}\ln |4x^2-4x+3| + \dfrac{1}{\sqrt{2}}\operatorname{Bgtan}\Lh \dfrac{2x-1}{\sqrt{2}}\Rh + c$\\


\begin{thebibliography}{999}
%\bibitem{a} Referentie
\end{thebibliography}

\end{document}